% Clase 12 --- La analogía de la manguera llena (§3.2).
% Resuelve la paradoja que el propio cálculo acababa de crear: si el electrón
% avanza 14 cm/hora, tardaría más de un día en llegar del interruptor a la
% lámpara, y sin embargo la luz se prende al instante. La clave es que la
% manguera YA está llena: lo que sale por la punta no es lo que entró.
\begin{center}
\begin{tikzpicture}[scale=1]

  % =============== (a) la manguera ===============
  \begin{scope}
    % la señal, inmediata
    \draw[figvecres, line width=1.2pt] (-2.6,1.15) -- (2.4,1.15);
    \node[figlbl, figred, anchor=south] at (-0.1,1.2)
      {la presión se transmite al instante};

    % canilla
    \draw[figgray, line width=1.2pt] (-3.3,-0.45) rectangle (-2.8,0.45);
    \node[figlblaux, anchor=east] at (-3.4,0) {canilla};

    % manguera, ya llena
    \fill[figblue!14] (-2.8,-0.32) rectangle (2.6,0.32);
    \draw[figgray, line width=1pt] (-2.8,-0.32) -- (2.6,-0.32);
    \draw[figgray, line width=1pt] (-2.8,0.32) -- (2.6,0.32);
    \node[figlblaux] at (0.1,0) {ya estaba llena de agua};

    % la molécula que entra y la que sale: distintas
    \node[figcarga, fill=figred, minimum size=6pt] at (-2.55,0) {};
    \node[figlbl, figred, anchor=south, scale=0.85] at (-2.5,0.42) {entra ésta};
    \node[figcarga, fill=figblue, minimum size=6pt] at (2.35,0) {};
    \node[figlbl, figblue, anchor=south, scale=0.85] at (2.3,0.42) {sale ésta};

    % avance lento
    \draw[figvecaux, figred, line width=0.9pt,
          -{Latex[length=2.2mm,width=1.6mm]}] (-2.5,-0.75) -- (-1.95,-0.75);
    \node[figlbl, figred, anchor=west, scale=0.9] at (-1.85,-0.75)
      {el agua avanza lento};

    \node[figlbl, anchor=north] at (-0.1,-1.35) {(a) manguera llena};
  \end{scope}

  % =============== (b) el circuito ===============
  \begin{scope}[yshift=-4.0cm]
    \draw[figvecres, line width=1.2pt] (-2.6,1.15) -- (2.4,1.15);
    \node[figlbl, figred, anchor=south] at (-0.1,1.2)
      {la señal eléctrica viaja rapidísimo};

    % interruptor
    \draw[figgray, line width=1.2pt] (-3.3,-0.45) rectangle (-2.8,0.45);
    \node[figlblaux, anchor=east] at (-3.4,0) {interruptor};

    % cable
    \fill[figblue!14] (-2.8,-0.32) rectangle (2.6,0.32);
    \draw[figgray, line width=1pt] (-2.8,-0.32) -- (2.6,-0.32);
    \draw[figgray, line width=1pt] (-2.8,0.32) -- (2.6,0.32);
    \node[figlblaux, scale=0.92] at (-0.15,0) {ya estaba lleno de electrones};

    % lámpara
    \draw[figamber, line width=1.2pt, fill=figamber!25] (3.0,0) circle (0.4);
    \foreach \a in {30,90,...,330} {
      \draw[figamber, line width=0.7pt] ($(3.0,0)+(\a:0.5)$) -- ($(3.0,0)+(\a:0.7)$);
    }
    \node[figlblaux, anchor=west, align=left] at (3.85,0)
      {enciende\\[-0.2em] al instante};

    \node[figcarga, fill=figred, minimum size=6pt] at (-2.55,0) {};
    \node[figcarga, fill=figblue, minimum size=6pt] at (2.35,0) {};

    \draw[figvecaux, figred, line width=0.9pt,
          -{Latex[length=2.2mm,width=1.6mm]}] (-2.5,-0.75) -- (-1.95,-0.75);
    \node[figlbl, figred, anchor=west, scale=0.9] at (-1.85,-0.75)
      {$v_d \sim$ algunos cm/hora};

    \node[figlbl, anchor=north] at (-0.1,-1.35) {(b) el mismo argumento en el cable};
  \end{scope}

\end{tikzpicture}\\[0.5em]
{\small\itshape La paradoja aparece al tomar el resultado del cálculo demasiado
literalmente: a algunos centímetros por hora, el electrón que sale del
interruptor tardaría más de un día en llegar a la lámpara, y sin embargo la luz
se prende ni bien se cierra el circuito. La salida es que el agua que sale por
la punta \textbf{no es} la que entró por la canilla: la manguera ya estaba llena
y todas las moléculas empiezan a moverse casi a la vez. En el cable pasa lo
mismo, con los electrones de conducción que ya estaban ahí. Conviene separar
tres velocidades que suelen confundirse: la \emph{real} de cada electrón
(enorme y errática), la de \emph{deriva} (lentísima, y la única que transporta
carga neta) y la de la \emph{señal} ---la que se percibe al encender la luz---,
que es muchísimo mayor que las otras dos.}
\end{center}
